% OLD TEXT FROM HERE ON
% This chapter contains text from the original submission.
% It will be revised as part of the salvage plan.

\section{CGE models with CES functions as an intermediate approach towards modeling production} \label{sec:ces}

Traditionally, many CGE models build on a Cobb-Douglas function to model consumption and production, where some nested applications also combine a Cobb-Douglas function for modeling substitution between intermediate goods and factors with a (nested) Leontief production function to capture constraints on the substitution among intermediates \citep[e.g.][]{GAMS.IO}. However, in the past decades the use of constant elasticity of substitution (CES) functions has become more common to model production as well as consumption in such frameworks \citep{Taylor1990,Klump.2012,Kim.2017}.

As indicated, these CES functions are used to model three types of trade-offs: while a CES utility function is employed to represent substitution between different consumption goods (captured by the substitution elasticity $\sigma$, see equation \eqref{eq:cesdemand}), a \textit{nested} CES production function can capture both, substitution between intermediates and factors (with elasticity $\theta$) as well as substitution between different intermediate goods in the nested part (with elasticity $\epsilon$) as shown in equation \eqref{eq:cesprodfun}.

An integral property of CES functions is the introduction of these explicit substitution parameters, which endow related models with additional flexibility. Specifically, these parameters provide a conceptual bridge between a Leontief production function and a Cobb-Douglas function: with low elasticities CES functions behave similar to the former and eventually converge to a Leontief function in the asymptotic limit (i.e., for $\varepsilon \to 0$), while an elasticity $\varepsilon$ equal to one coincides with the Cobb-Douglas case.

Hence, in Kuhnian terms such CES models are more general insofar as they capture not only both archetypical representations of production processes associated with different paradigmatic perspectives, but also the relevant middle ground or paradigmatic gap, that arises from competing assumptions. In this context, the Leontief production function is interpreted as a heterodox archetype for modeling production in the short run that originates in the works of classical economists like Smith or Ricardo \citep[e.g.][]{Blecker2019}, whereas the Cobb-Douglas case is taken to represent the archetypical neoclassical approach. This interpretation also follows from section~\ref{sec:IO_models}, where the choice of production function was identified as a key paradigmatic wedge. Hence, an approach based on CES-functions allows for exploring the consequences of this major conceptual difference.

In what follows we compare the effect of a sectoral shock to final demand as defined in section \ref{sec:data} across different modeling strategies and specifications. The aim of this exercise is to better assess how the choice of model translates into different predictions of the economic effects of transformational policies building on increased public investment. In our view, this aspect is a key factor for assessing whether and to what extent the use of CES functions may contribute to bridging paradigmatic gaps that separate different modeling traditions.\footnote{Note that to initialize the demand shock formally, we can simply enter additional investment as gross sums in the standard Leontief model as indicated in equation \eqref{eq:leontief_effect}. As already indicated in section \ref{sec:data}, for the general equilibrium model, which is formulated in normalized terms, changes in final demands have to be entered relative to current final spending in affected sectors.}

\subsection{Basic properties of CGE models with (nested) CES functions} \label{subsec:introducingCES}

In line with other related applications, \citet{BF2019, BF2022} employ CES functions in various contexts. In what follows we focus on the treatment as presented in \citet{BF2019}, which not only makes consistent use of CES-functions, but also suggests specific values for the elasticities of substitution based on empirical considerations, namely $(\varepsilon,\theta,\sigma) = (0.001,0.5,0.9)$. Thereby, the low elasticity between intermediate goods is of particular interest as this choice implies that the CES production function mimics a Leontief production function in the nested part of equation \eqref{eq:cesprodfun} shown below. Although this can be interpreted as taking one step towards a more heterodox perspective on modeling production, the full implications of this modification are not ex ante clear, however, also because the final results obtained also depend on the other elasticities employed in a specific application.

\citet{BF2019} employ the following constant returns to scale CES function for modeling final demand, which aligns closely with the exposition in section \ref{subsec:CGE}, especially equation \eqref{eq:cgutilitymax}.

\begin{equation}
    \label{eq:cesdemand}
    \mathcal D(\mathbf c)= \left( \sum_{i=0}^N \omega_{i0}^{\frac 1 \sigma} c_i^{\frac {\sigma-1}{\sigma}} \right)^{\frac {\sigma}{\sigma-1}}
\end{equation}

For the production side the authors also retain the assumption of constant returns to scale. As indicated before and in line with the standard use of nested CES functions they assume specific elasticities for modeling substitution across intermediate goods ($\epsilon$) as well as substitution between production factors and intermediates ($\theta$). 

Doing so they arrive at the formulation

\begin{equation}
    \label{eq:cesprodfun}
    f_u(l_u, \mathbf x_{\mathbf{s} u}) =  \mathcal A_u \left(\omega_{lu}^{\frac{1}{\theta}} l_u^{\frac {\theta-1}{\theta}} + (1- \omega_{lu})^{\frac{1}{\theta}} \left( \left(\sum_{k=1}^N \omega_{ku}^{\frac{1}{\epsilon}} x_{ku}^{\frac{\varepsilon - 1}{\varepsilon}} \right)^{\frac{\varepsilon}{\varepsilon - 1}} \right)^{\frac{\theta-1}{\theta}} \right)^{\frac{\theta}{\theta - 1}}.
\end{equation}

Taken together with the market clearing conditions for goods markets (equation \eqref{eq:_goods_clear}) and the labor market clearing conditions as employed in problems \eqref{eq:profitmaxproblem} and \eqref{eq:costmin} these formulations allow for the derivation of new equilibrium configurations obtained after imposing specific shocks.\footnote{While most standard textbook representations of the CES-function do not include the substitution elasticities as exponents of the share parameters, ours does. Intuitively, incorporating the substitution elasticities in that way reflects a possible adjustment process, because it not only allows for the substitution of inputs, but also for the change in effective contribution of each input. Adding this exponent is in line with most practical applications of the model \citep[e.g.][]{Kim.2017}, including the implementation of the model in \cite{BF2019}, which, in this instance, deviates from the equations presented in the paper.} When doing so, combinations of efficient prices and quantities have to be calculated iteratively until a stable solution is attained, which then represents the new equilibrium.\footnote{In standard textbooks \citep[e.g.][]{cardeneteAppliedGeneralEquilibrium2017} the existence of equilibria is usually motivated by convergence theorems of fixed-point iterations like Bowers' fixed point theorem. In practice however, nonlinear root finding methods like modifications of Newton's method are employed to find new equilibria. In our implementation we used the Nonlinearsolve.jl package \citep[][]{palNonlinearSolvejlHighPerformanceRobust2025} implemented in Julia \citep{bezansonJuliaFreshApproach2017}. There also exist generalized methods, where the cost function and demand function do not have to be specified algebraically \citep[e.g.][]{choiComplementarityApproachSolving2014}.} When substituting the concrete expressions for the cost function, the wage function, and the demand function into equation~\eqref{eq:equilibrium}, projections for efficient prices and quantities are provided by the following pair of equations for each $1 \leq u \leq n$.

\begin{gather}
\begin{split}    
  p_u = \biggl( \mathcal A_u^{\varepsilon - 1} \left( \omega_{ul} w_u ^ {1 - \varepsilon} + (1 - \omega_{ul})  q_u ^ {1 - \varepsilon} \right) \biggr) ^ {\frac{1}{1 - \varepsilon}} \\
  y_u = p_u^{-\theta} \left(\sum_{s=1}^n \omega_{su}  p_s ^ {\varepsilon}  \mathcal A_s ^ {\varepsilon - 1} q_s ^ {\theta - \varepsilon}  (1 - \omega_{ls} )  y_s \right) +  c_u,
  \end{split}
\end{gather}
where 
\begin{gather}
\begin{split}
  q_u = \left( \sum_{s=1}^n \omega_{us}  p_s ^ {1 - \theta} \right) ^ {\frac{1}{1 - \theta}}. \\
  w_u = p_u \mathcal A_u ^ {\frac {\varepsilon - 1}{\varepsilon}} \omega_{lu} ^ {\frac{1}{\varepsilon}} y_u ^ {\frac{1}{\varepsilon}} l_u ^ {\frac {-1}{\varepsilon}} \\
  c_u = D_u \left( \frac{p_u}{\text{CPI}} \right)^ {-\theta} \omega_{0u} \sum_{s=1}^n w_s l_s. \\ CPI = \left( \sum_{u=1}^n \omega_{0uu} p_u^{\theta  -1} \right)^{\frac{1}{1\theta - 1}}.
  \end{split}
\end{gather}

Note that, in this formulation $\mathbf{q}$ represent the prices of intermediary goods. Taken together with equation~\eqref{eq:market_clearing_sheppard}, we observe from the above system of equations that new equilibrium prices are given by a fixed point of the cost function at which the new quantities satisfy the market clearing conditions. Figuratively speaking this implies that the above equations represent an iterative algorithm, which mimics the intuition underlying the Walrasian auctioneer, who successively adjusts prices until an efficient allocation is reached \citep{cardeneteAppliedGeneralEquilibrium2017}.

Following the descriptions in section \ref{subsec:CGE} post shock relative prices are obtained by the normalization with the chosen numeraire and increase in real GDP is calculated with the Laspeyres index, via the consumption side.  Nominal GDP, which is of little interest to us in this application can be either obtained from the demand side via $\sum_{s = 1}^n \post {\tilde p_s }\post c_s$ or from the supply side via $\sum_{s = 1}^n \post {\tilde w_s } l_s$. Moreover the sectoral composition of intermediate or final demand can be assessed and compared to the pre-shock equilibrium by comparing the Domar weights and consumption shares.

 \subsection{To what extent do elasticities matter?}\label{subsec:CGE_elasticity}

To assess the implications of this setup and the variation in outcomes associated with changing assumptions on substitution elasticities, we compare the estimated reaction of real GDP to a transformational demand shock as specified in section \ref{sec:data} for a variety of scenarios. Specifically we vary individual elasticities between 0 and 1, while holding all other elasticities constant at some pre-specified level to explore the full range of possible outcomes associated with this specific axiomatic variation.

In this vein, the panel in Figure \ref{fig:elasticity-gradient} shows estimated GDP for varying elasticities of all types assuming that all other elasticities are held constant at $0.9$ (upper left corner), $0.5$ (upper right corner), $0.2$ (bottom left corner) and $0.1$ (bottom right corner). For purposes of comparison these graphical representations also include estimates taken from a traditional Leontief model (as explained in section \ref{sec:leontief}) as well as a standard Cobb Douglas approach  to assess what difference the introduction of more flexible and heterogenous assumptions on substitutability properties eventually brings. 

\begin{figure}
    \centering
    \includegraphics[width=1\linewidth]{pictures/impulse.png}
    \caption{Effect of different elasticities on real GDP}
    \label{fig:elasticity-gradient}
\end{figure}

An overall inspection of the results show that in this comparison the standard Leontief approach delivers the largest GDP-estimates -- on this basis GDP is expected to grow by about 1.5\%, which is equivalent to a multiplier of 1.15 given that the initial investment impulse was specified to be 40.337 bn. €, about 1.3\% of GDP. In contrast, the Cobb-Douglas specification predicts a small reduction in GDP that results from efficiency-losses arising from the increase in green investment. These losses are due to the fact that the imposed green investment leads to a forced reallocation of intermediates and production factors towards the expanding sectors, which alters an already optimal allocation and, hence, induces inefficiencies. In turn, real output decreases slightly and the major macroeconomic impact of a green investment demand shock in this framework is to induce sector-specific inflationary pressures.

With regard to how changes in assumptions on substitution elasticities impact on final outcomes Figure \ref{fig:elasticity-gradient} indicates that real GDP is typically increasing in all three elasticities employed -- the higher the chosen elasticities, the higher are the respective GDP-estimates. This finding reproduces the intuition that easier substitution will lead to higher (aggregate) output and, hence, CES-based estimates are consistently lower than their Cobb-Douglas counterparts.

Against this backdrop we concede that the introduction of Leontief-like properties in the production function helps bridging paradigmatic cleavages on a level of conceptual foundations -- e.g. with respect to the question how to adequately represent production processes. However, we find that introducing such a conceptual shift in neoclassical IO-models will actually serve to increase the gap between both approaches in terms of predicted outcomes. This result is not too surprising when taking into account that the assumption of a Leontief technology implies less flexibility as compared to a Cobb-Douglas technology. Hence, in the context of assumptions on substitutability neoclassical approaches are more optimistic than the heterodox archetype: viewed in isolation greater substitutability comes with higher prospects for growth in case of demand shocks and, correspondingly, leads to a more positive assessment of fiscal intervention in terms of forced green investment. However, this relatively more optimistic stance does in no way compensate for the fact that the estimated growth effects remain negative, which is effectively driven by assumptions on capacity constraints. However, before turning to this key aspect in section \ref{sec:labor} further below, we first inspect structural differences between model outcomes on a more granular level.

\subsection{The effect of a demand-shock in inter-paradigmatic comparison?}

To facilitate a comparative discussion on characteristic patterns associated with the imposition of exogenous demand shocks we use the opportunity to illuminate some structural differences in model responses between the two types of applications.

First and foremost, we note that in the standard Leontief model the initial change in final demand expenditures is, somewhat unsurprisingly, fully reflected by an expansion of final demand $\mathbf{c}$ with all prices staying constant.\footnote{On top of the initial stimulus, the Leontief model also features an induced demand effect arising from the growth in income resulting from the original expansion. However, these effects are too small to be visible in figure~\ref{fig:comparison-sec4}.} This pure quantity expansion reflects the output adjustment logic employed in the underlying model closure and stands in contrast to the more heterogeneous results provided by the CGE approach. In this setup an exogenous demand shock is accommodated by adjusting total revenues to the (net) change in monetary expenditure, which can happen through an expansion of either quantities or prices. And, indeed, in those sectors in which the exogenous shock hits, (relative) prices show a substantial increase and contribute significantly to match the demand stimulus. Quantities, on the other hand, only rise in five out of seven shocked sectors as is indicated by the left panel in Figure \ref{fig:comparison-sec4}. Moreover, in all shocked sectors the combined changes of prices and quantities do not suffice to fully reflect the imposed expansion in expenditures so that a third component is required to eventually accommodate the external stimulus. This third component takes the form of a crowding-out effect, i.e., a decrease of final consumption net of the additional stimulus. These sectoral crowding-out effects, which appear mainly in sectors primarily affected from the exogenous shock, is complemented by aggregate crowding-out effect driven by the overall scarcity of production factors ($l^{max}$ in our model setup), which implies that even sectors relatively unaffected from the exogenous shock in terms of their own production or intermediate good relationships will show a slight decrease in gross output (as shown in Figure \ref{fig:sectoral-changes} below).

\begin{figure}
  \begin{center}
    \includegraphics[width=0.95\textwidth]{pictures/panel.png}
  \end{center}
  \caption{Changes in prices and consumption for shocked sectors (left panel) and quantities and prices for all sectors (right panel). Outputs are based on a model run using the standard specification from \citet{BF2019} with $(\varepsilon,\theta,\sigma) = (0.001, 0.5, 0.9)$. Error bars indicate the range of outcomes that can be achieved for extreme parameter values (specifically, $(\varepsilon,\theta,\sigma) = (0.001, 0.1, 0.6)$ and $(\varepsilon,\theta,\sigma) = (0.99, 0.99, 0.99)$).}
  \label{fig:comparison-sec4}
\end{figure} 

In general, we observe the theoretically expected positive correlation between quantity and price changes in the response of the CGE model as shown in the right panel of Figure \ref{fig:comparison-sec4}. This pattern is consistent with rationalizing the exogenous shock in direct analogy to positive demand shocks as employed in standard microeconomics, which shift demand curves outwards. Hence, relative prices of those goods, that are in high demand, increase, while prices for the remaining goods fall, delivering a correlation that runs counter the usual notion of the law of demand\footnote{Shifting a curve invokes an exception from the standard law of demand, which effectively represents a recourse on a ceteris paribus clause, see \citet{Hausman.1990}}. In our application this pattern holds not only for five out of seven sectors subject to the exogenous demand shock (colored as red), but also for the large majority of remaining sectors (shown in pink). Important outliers to this overall development are five specific sectors appearing in the second quadrant in the right panel of Figure \ref{fig:comparison-sec4}, where two belong to the shocked sectors (again in red), whereas three others are not directly affected from the shock.\footnote{The five sectors in the second quadrant are \textit{Lumber and Wood}, \textit{Metal ores and Mining} and \textit{Coke and refined petroleum products} as well as \textit{Chemicals and chemical products} and \textit{Electrical equipment} (shocked; in red).} Those sectors instead mimic the law of demand when showing rising prices and decreasing quantities. Also, these sectors show a stronger average decline in terms of quantity then the remaining sectors, that are not subject to the exogenous shock. This pattern occurs because of a 'cost-shock', where higher prices for intermediate goods originating in the shocked sectors translate into a net price increase for these five sectors that coins aggregate outcomes. The very same reason explains, why these sectors show an over-proportional reduction in quantities as they are not only subject to the aggregate crowding out effects (due to scarcity of labor), but also more sensitive to the transmission of sectoral crowding out effects occurring in shocked sector (which leads to shortage of intermediate goods). 

\begin{figure}
    \centering
    \includegraphics[width=1\linewidth]{pictures/diff_lambda_imp.png}
    \caption{Change in outputs, based on a model run using the standard specification  from \citet{BF2019}   with $(\varepsilon,\theta,\sigma) = (0.001, 0.5, 0.9)$. Shocked sectors remain unshaded.}
    \label{fig:sectoral-changes}
\end{figure}

Finally, by assessing changes in sectoral production after the demand shock in Figure \ref{fig:sectoral-changes}, we find that the output adjustment approach will lead to an expansion of almost all sectors, while the general equilibrium approach leads to a reduction of gross output in the large majority of sectors, that are not directly affected by the demand shock. Moreover, the shock also depresses gross output and final consumption associated with two of the shocked sectors (in line with Figure \ref{fig:comparison-sec4} above, which focuses on final demand only). This reproduces the finding from section \ref{subsec:CGE_elasticity} that a demand shock will depress real GDP on a sectoral level as the (comparatively lower) average expansion of shocked sectors does under no conditions fully compensate for the reduction in the production of unaffected sectors. Notwithstanding this overall pattern, the complex substitution dynamics implied the CGE approach lead to exceptional outcomes in specific sectors. The most obvious idiosyncratic outcome occurs in two sectors that are not shocked –– \textit{building construction works (sector 33 )} and \textit{civil engineering works (sector 34)} --, but, rather, massively expand their intermediate goods production at the expense of final outputs, which represents an over-proportional intersectoral crowding-out effect. In other words, intermediary goods from these sectors are in extremely high demand due to the enforced shift in final expenditures. For both of these sectors the expansion in intermediate goods production is so strong, that these sectors expand more strongly than in the Leontief case and thereby partially compensate for the lack of freely available workers in shocked sectors.

As already indicated, we observe that the change in assumptions with regard to production technology, that does represent a step towards a more heterodox perspective, does not contribute to closing the paradigmatic gap in terms of results. As this exemplary analysis shows, the neoclassical outcome is mostly characterized by the generally constraining scarcity of factors that rationalizes any directed exogenous demand shock as an additional burden, requiring the potentially costly or inefficient reallocation of scarce production factors. By moving towards a more Leontief-like understanding of production processes such adaption and reallocation actually becomes harder to accomplish. Hence, this axiomatic variation in isolation has the curious effect to increase paradigmatic distance in terms results, although it reduces the same distance on the level of assumptions.

